\section*{Abstract}

This paper addresses the problem of \emph{structural diagnosability} of enterprise states under conditions of incomplete and indirect observability.
Observable metrics, indicators, and outcome variables generally fail to uniquely determine the internal structural configurations that generate them.
As a result, distinct enterprise structures may induce indistinguishable observable projections, imposing intrinsic limits on diagnostic inference.

Within the formal framework of the Enterprise Continuum \(K_{EA}\) and its executable specification ETS.K\_EA.v1.1, structural diagnosability is analyzed as an epistemic property of the mapping from internal structural spaces to observable domains.
The analysis is strictly diagnostic and non-prescriptive: it does not address control, optimization, intervention, or corrective action.
Its purpose is to characterize when an enterprise state can be inferred to belong to a \emph{structural phase class}, rather than identified as a unique configuration.

Three contributions are developed.
First, structural diagnosability is formally defined in terms of distinguishability between equivalence classes induced by observable projections.
Second, diagnostic claims are interpreted within a phase-based framework, establishing that, in general, well-founded diagnoses are phase-level rather than state-level assertions.
Third, falsifiability conditions for diagnostic statements are specified by identifying observable distinctions whose occurrence would refute phase membership.

The scope of the results is intentionally limited.
No claims are made regarding desirable states, optimal structures, or recommended actions.
Non-diagnosability is treated as a valid and informative outcome, reflecting fundamental epistemic constraints rather than deficiencies of the model.
Under these conditions, structural diagnosability is established as a formally bounded property of enterprise models, jointly determined by observability and structural organization.
